\documentclass[12pt]{article}
\usepackage{latexblog}

% ============================================================
% Blog Metadata
% ============================================================
\blogtitle{破解Db Schema序列号}
\blogdate{2016-07-22}
\blogtags{刨根问底, 探索分析}
\bloglang{zh}
\blogsummary{}

\begin{document}
\makeblogtitle

我始终认为数据库设计在系统设计中是一个很重要的工作，然而一直没有比较好的ER建模工具。使用过MySQL Workbench和Power Designer两种工具，但都存在很多不喜欢的地方，直到遇到DbSchema后眼前一亮，这才是一个Nice的工具嘛。 很可惜对于我们这种屌丝来说，是不舍得花钱去购买一个license的，试用期15天到了怎么办呢？当时也没发现有可用的破解版，因为它是基于Java的，这对破解来说减小了难度，于是趁着辞职后在家没事的空档来研究了一下破解。其实也就上午花了一会时间就搞定了。记录下破解的过程。

首先是找到dbschema.jar，这是程序的主要jar包，其他是一些第三方的jar包和jdbc驱动等，于是它就是破解的关键。利用jd-gui反编译这个jar包，首先把源码都保存下来。

顺藤摸瓜，首先打开dbschema的注册窗口，根据里面的关键字搜索，比如Registration，然后一个个去找，这时，发现一个对话框：

\begin{Shaded}
\begin{Highlighting}[]
\KeywordTok{public} \KeywordTok{class}\NormalTok{ RegistrationDialog}
\KeywordTok{...}
\BuiltInTok{JButton}\NormalTok{ localJButton1 }\OperatorTok{=} \KeywordTok{new} \BuiltInTok{JButton}\OperatorTok{(}\FunctionTok{getAction}\OperatorTok{(}\StringTok{"register"}\OperatorTok{));}
\end{Highlighting}
\end{Shaded}

这不就是注册的按钮么？然后就看它的action:

\begin{Shaded}
\begin{Highlighting}[]
\CommentTok{/*     */}   \KeywordTok{public} \DataTypeTok{void} \FunctionTok{register}\OperatorTok{()} \OperatorTok{\{}
\CommentTok{/*  96 */}     \BuiltInTok{String}\NormalTok{ str1 }\OperatorTok{=} \KeywordTok{this}\OperatorTok{.}\FunctionTok{b}\OperatorTok{.}\FunctionTok{getText}\OperatorTok{();}
\CommentTok{/*     */}     
\CommentTok{/*  98 */}     \ControlFlowTok{if} \OperatorTok{((}\NormalTok{str1 }\OperatorTok{==} \KeywordTok{null}\OperatorTok{)} \OperatorTok{||} \OperatorTok{(}\NormalTok{str1}\OperatorTok{.}\FunctionTok{length}\OperatorTok{()} \OperatorTok{==} \DecValTok{0}\OperatorTok{))} \OperatorTok{\{}
\CommentTok{/*  99 */}       \BuiltInTok{JOptionPane}\OperatorTok{.}\FunctionTok{showMessageDialog}\OperatorTok{(}\KeywordTok{this}\OperatorTok{,}\NormalTok{ d}\OperatorTok{.}\FunctionTok{a}\OperatorTok{(}\DecValTok{11}\OperatorTok{),} \StringTok{"Error"}\OperatorTok{,} \DecValTok{0}\OperatorTok{);}
\CommentTok{/* 100 */}       \ControlFlowTok{return}\OperatorTok{;}
\CommentTok{/*     */}     \OperatorTok{\}}
\CommentTok{/* 102 */}     \BuiltInTok{String}\NormalTok{ str2 }\OperatorTok{=} \KeywordTok{this}\OperatorTok{.}\FunctionTok{c}\OperatorTok{.}\FunctionTok{getText}\OperatorTok{();}
\CommentTok{/* 103 */}     \ControlFlowTok{if} \OperatorTok{(}\NormalTok{str2 }\OperatorTok{==} \KeywordTok{null}\OperatorTok{)} \OperatorTok{\{}
\CommentTok{/* 104 */}       \BuiltInTok{JOptionPane}\OperatorTok{.}\FunctionTok{showMessageDialog}\OperatorTok{(}\KeywordTok{this}\OperatorTok{,}\NormalTok{ d}\OperatorTok{.}\FunctionTok{a}\OperatorTok{(}\DecValTok{19}\OperatorTok{),} \StringTok{"Error"}\OperatorTok{,} \DecValTok{0}\OperatorTok{);}
\CommentTok{/* 105 */}       \ControlFlowTok{return}\OperatorTok{;}
\CommentTok{/*     */}     \OperatorTok{\}}
\CommentTok{/* 107 */}\NormalTok{     str1 }\OperatorTok{=}\NormalTok{ str1}\OperatorTok{.}\FunctionTok{trim}\OperatorTok{();}
\CommentTok{/* 108 */}\NormalTok{     str2 }\OperatorTok{=}\NormalTok{ str2}\OperatorTok{.}\FunctionTok{trim}\OperatorTok{();}
\CommentTok{/* 109 */}\NormalTok{     e}\OperatorTok{.}\FunctionTok{b}\OperatorTok{(}\NormalTok{d}\OperatorTok{.}\FunctionTok{a}\OperatorTok{(}\DecValTok{31}\OperatorTok{),}\NormalTok{ str1}\OperatorTok{);}
\CommentTok{/* 110 */}\NormalTok{     e}\OperatorTok{.}\FunctionTok{b}\OperatorTok{(}\NormalTok{d}\OperatorTok{.}\FunctionTok{a}\OperatorTok{(}\DecValTok{21}\OperatorTok{),}\NormalTok{ str2}\OperatorTok{);}
\CommentTok{/*     */}     
\CommentTok{/* 112 */}     \DataTypeTok{int}\NormalTok{ i }\OperatorTok{=}\NormalTok{ g}\OperatorTok{.}\FunctionTok{b}\OperatorTok{();}
\CommentTok{/* 113 */}     \ControlFlowTok{if} \OperatorTok{(}\NormalTok{i }\OperatorTok{==} \BuiltInTok{Integer}\OperatorTok{.}\FunctionTok{MAX\_VALUE}\OperatorTok{)} \OperatorTok{\{}
\CommentTok{/* 114 */}       \FunctionTok{dispose}\OperatorTok{();}
\CommentTok{/* 115 */}       \BuiltInTok{JOptionPane}\OperatorTok{.}\FunctionTok{showMessageDialog}\OperatorTok{(}\KeywordTok{this}\OperatorTok{.}\FunctionTok{a}\OperatorTok{.}\FunctionTok{c}\OperatorTok{(),}\NormalTok{ d}\OperatorTok{.}\FunctionTok{a}\OperatorTok{(}\DecValTok{23}\OperatorTok{),} \StringTok{"Info"}\OperatorTok{,} \DecValTok{1}\OperatorTok{,} \KeywordTok{null}\OperatorTok{);}
\CommentTok{/* 116 */}       \KeywordTok{this}\OperatorTok{.}\FunctionTok{a}\OperatorTok{.}\FunctionTok{c}\OperatorTok{().}\FunctionTok{c}\OperatorTok{();}
\CommentTok{/* 117 */}     \OperatorTok{\}} \ControlFlowTok{else} \ControlFlowTok{if} \OperatorTok{((}\NormalTok{i }\OperatorTok{\textgreater{}} \DecValTok{0}\OperatorTok{)} \OperatorTok{\&\&} \OperatorTok{(}\NormalTok{str1}\OperatorTok{.}\FunctionTok{toLowerCase}\OperatorTok{().}\FunctionTok{startsWith}\OperatorTok{(}\StringTok{"extend"}\OperatorTok{)))} \OperatorTok{\{}
\CommentTok{/* 118 */}       \FunctionTok{dispose}\OperatorTok{();}
\CommentTok{/* 119 */}       \BuiltInTok{JOptionPane}\OperatorTok{.}\FunctionTok{showMessageDialog}\OperatorTok{(}\KeywordTok{this}\OperatorTok{.}\FunctionTok{a}\OperatorTok{.}\FunctionTok{c}\OperatorTok{(),}\NormalTok{ d}\OperatorTok{.}\FunctionTok{a}\OperatorTok{(}\DecValTok{24}\OperatorTok{).}\FunctionTok{replaceAll}\OperatorTok{(}\StringTok{"}\SpecialCharTok{\textbackslash{}\textbackslash{}}\StringTok{\{days}\SpecialCharTok{\textbackslash{}\textbackslash{}}\StringTok{\}"}\OperatorTok{,} \StringTok{""} \OperatorTok{+}\NormalTok{ i}\OperatorTok{),} \StringTok{"Info"}\OperatorTok{,} \DecValTok{1}\OperatorTok{,} \KeywordTok{null}\OperatorTok{);}
\CommentTok{/* 120 */}     \OperatorTok{\}} \ControlFlowTok{else} \ControlFlowTok{if} \OperatorTok{(}\NormalTok{i }\OperatorTok{==} \OperatorTok{{-}}\DecValTok{2}\OperatorTok{)} \OperatorTok{\{}
\CommentTok{/* 121 */}       \BuiltInTok{String}\NormalTok{ str3 }\OperatorTok{=}\NormalTok{ d}\OperatorTok{.}\FunctionTok{a}\OperatorTok{(}\DecValTok{77}\OperatorTok{).}\FunctionTok{replace}\OperatorTok{(}\StringTok{"\{0\}"}\OperatorTok{,} \KeywordTok{new} \BuiltInTok{SimpleDateFormat}\OperatorTok{(}\StringTok{"dd.MMMMM.yyyy"}\OperatorTok{).}\FunctionTok{format}\OperatorTok{(}\KeywordTok{new} \BuiltInTok{Date}\OperatorTok{(}\NormalTok{g}\OperatorTok{.}\FunctionTok{c}\OperatorTok{())));}
\CommentTok{/* 122 */}       \BuiltInTok{JOptionPane}\OperatorTok{.}\FunctionTok{showMessageDialog}\OperatorTok{(}\KeywordTok{this}\OperatorTok{.}\FunctionTok{a}\OperatorTok{.}\FunctionTok{c}\OperatorTok{(),}\NormalTok{ str3}\OperatorTok{,} \StringTok{"Error"}\OperatorTok{,} \DecValTok{0}\OperatorTok{);}
\CommentTok{/*     */}     \OperatorTok{\}} \ControlFlowTok{else} \OperatorTok{\{}
\CommentTok{/* 124 */}       \BuiltInTok{JOptionPane}\OperatorTok{.}\FunctionTok{showMessageDialog}\OperatorTok{(}\KeywordTok{this}\OperatorTok{.}\FunctionTok{a}\OperatorTok{.}\FunctionTok{c}\OperatorTok{(),}\NormalTok{ d}\OperatorTok{.}\FunctionTok{a}\OperatorTok{(}\DecValTok{8}\OperatorTok{),} \StringTok{"Error"}\OperatorTok{,} \DecValTok{0}\OperatorTok{);}
\CommentTok{/*     */}     \OperatorTok{\}}
\CommentTok{/*     */}   \OperatorTok{\}}
\end{Highlighting}
\end{Shaded}

112行开始有点意思，其实大概能猜到是干什么，反正是算剩余天数的，那么这个int i = g.b();就是最核心的东西了：

\begin{Shaded}
\begin{Highlighting}[]
\KeywordTok{public} \DataTypeTok{static} \DataTypeTok{int} \FunctionTok{b}\OperatorTok{()}
  \OperatorTok{\{}
    \BuiltInTok{String}\NormalTok{ str1 }\OperatorTok{=}\NormalTok{ e}\OperatorTok{.}\FunctionTok{d}\OperatorTok{(}\NormalTok{d}\OperatorTok{.}\FunctionTok{a}\OperatorTok{(}\DecValTok{31}\OperatorTok{),} \KeywordTok{null}\OperatorTok{);}
    \BuiltInTok{String}\NormalTok{ str2 }\OperatorTok{=}\NormalTok{ e}\OperatorTok{.}\FunctionTok{d}\OperatorTok{(}\NormalTok{d}\OperatorTok{.}\FunctionTok{a}\OperatorTok{(}\DecValTok{21}\OperatorTok{),} \KeywordTok{null}\OperatorTok{);}
    \DataTypeTok{int}\NormalTok{ m }\OperatorTok{=} \OperatorTok{{-}}\DecValTok{1}\OperatorTok{;}
    \ControlFlowTok{if} \OperatorTok{((}\NormalTok{str1 }\OperatorTok{!=} \KeywordTok{null}\OperatorTok{)} \OperatorTok{\&\&} \OperatorTok{(}\NormalTok{str2 }\OperatorTok{!=} \KeywordTok{null}\OperatorTok{)} \OperatorTok{\&\&} \OperatorTok{(}\NormalTok{str2}\OperatorTok{.}\FunctionTok{length}\OperatorTok{()} \OperatorTok{\textgreater{}} \DecValTok{3}\OperatorTok{))}
    \OperatorTok{\{}
      \ControlFlowTok{if} \OperatorTok{((}\NormalTok{str1}\OperatorTok{.}\FunctionTok{toLowerCase}\OperatorTok{().}\FunctionTok{startsWith}\OperatorTok{(}\StringTok{"extend"}\OperatorTok{))} \OperatorTok{\&\&} \OperatorTok{(}\FunctionTok{c}\OperatorTok{(}\NormalTok{str1}\OperatorTok{,}\NormalTok{ str2}\OperatorTok{)))}
      \OperatorTok{\{}
\NormalTok{        m }\OperatorTok{=} \BuiltInTok{Math}\OperatorTok{.}\FunctionTok{max}\OperatorTok{(}\DecValTok{15} \OperatorTok{{-}} \FunctionTok{f}\OperatorTok{(}\StringTok{"mmax"}\OperatorTok{),} \OperatorTok{{-}}\DecValTok{1}\OperatorTok{);}
      \OperatorTok{\}}
      \ControlFlowTok{else} \ControlFlowTok{if} \OperatorTok{(}\NormalTok{str2}\OperatorTok{.}\FunctionTok{length}\OperatorTok{()} \OperatorTok{\textgreater{}} \DecValTok{15}\OperatorTok{)}
      \OperatorTok{\{}
        \BuiltInTok{String}\NormalTok{ str3 }\OperatorTok{=}\NormalTok{ str2}\OperatorTok{.}\FunctionTok{substring}\OperatorTok{(}\DecValTok{4}\OperatorTok{,} \DecValTok{9}\OperatorTok{);}
        \BuiltInTok{String}\NormalTok{ str4 }\OperatorTok{=}\NormalTok{ str2}\OperatorTok{.}\FunctionTok{substring}\OperatorTok{(}\DecValTok{0}\OperatorTok{,} \DecValTok{4}\OperatorTok{)} \OperatorTok{+}\NormalTok{ str2}\OperatorTok{.}\FunctionTok{substring}\OperatorTok{(}\DecValTok{9}\OperatorTok{);}
        \ControlFlowTok{if} \OperatorTok{(}\FunctionTok{c}\OperatorTok{(}\StringTok{"ax5"} \OperatorTok{+}\NormalTok{ str1 }\OperatorTok{+} \StringTok{"b52w"} \OperatorTok{+}\NormalTok{ str3 }\OperatorTok{+} \StringTok{"vb3"}\OperatorTok{,}\NormalTok{ str4}\OperatorTok{))}
        \OperatorTok{\{}
          \ControlFlowTok{try}
          \OperatorTok{\{}
\NormalTok{            k }\OperatorTok{=} \BuiltInTok{Integer}\OperatorTok{.}\FunctionTok{parseInt}\OperatorTok{(}\NormalTok{str3}\OperatorTok{)} \OperatorTok{*} \DecValTok{86400000L} \OperatorTok{+} \DecValTok{1356994800000L}\OperatorTok{;}
          \OperatorTok{\}}
          \ControlFlowTok{catch} \OperatorTok{(}\BuiltInTok{NumberFormatException}\NormalTok{ localNumberFormatException}\OperatorTok{)}
          \OperatorTok{\{}
\NormalTok{            c}\OperatorTok{.}\FunctionTok{b}\OperatorTok{(}\NormalTok{localNumberFormatException}\OperatorTok{);}
          \OperatorTok{\}}
\NormalTok{          m }\OperatorTok{=} \BuiltInTok{Integer}\OperatorTok{.}\FunctionTok{MAX\_VALUE}\OperatorTok{;}
        \OperatorTok{\}}
      \OperatorTok{\}}
    \OperatorTok{\}}
    \ControlFlowTok{else} \OperatorTok{\{}
\NormalTok{      m }\OperatorTok{=} \BuiltInTok{Math}\OperatorTok{.}\FunctionTok{max}\OperatorTok{(}\DecValTok{15} \OperatorTok{{-}} \FunctionTok{f}\OperatorTok{(}\StringTok{"mma"}\OperatorTok{),} \OperatorTok{{-}}\DecValTok{1}\OperatorTok{);}
    \OperatorTok{\}}
    \ControlFlowTok{return}\NormalTok{ m}\OperatorTok{;}
  \OperatorTok{\}}
\end{Highlighting}
\end{Shaded}

看到这，我们其实已经拿到了计算key的方法，只不过这是一个验证的函数，如果我们要计算出key，需要反向推倒出来，这里就不具体解释了，最终反向出来的代码其实很简单，我做了一个C++版本的：

\begin{Shaded}
\begin{Highlighting}[]
\NormalTok{inline }\DataTypeTok{const}\NormalTok{ string }\FunctionTok{generateKey}\OperatorTok{(}\NormalTok{string name}\OperatorTok{)}
\OperatorTok{\{}
\NormalTok{    string salt }\OperatorTok{=} \FunctionTok{getSalt}\OperatorTok{();}
\NormalTok{    cout }\OperatorTok{\textless{}\textless{}} \StringTok{"salt:"} \OperatorTok{\textless{}\textless{}}\NormalTok{ salt }\OperatorTok{\textless{}\textless{}}\NormalTok{ endl}\OperatorTok{;}
\NormalTok{    string encryptSource }\OperatorTok{=} \StringTok{"ax5"} \OperatorTok{+}\NormalTok{ name }\OperatorTok{+} \StringTok{"b52w"} \OperatorTok{+}\NormalTok{ salt }\OperatorTok{+} \StringTok{"vb3"}\OperatorTok{;}
\NormalTok{    cout }\OperatorTok{\textless{}\textless{}} \StringTok{"encrypt:"} \OperatorTok{\textless{}\textless{}}\NormalTok{ encryptSource }\OperatorTok{\textless{}\textless{}}\NormalTok{ endl}\OperatorTok{;}
\NormalTok{    string hash }\OperatorTok{=} \FunctionTok{MD5}\OperatorTok{(}\NormalTok{encryptSource}\OperatorTok{).}\FunctionTok{toStr}\OperatorTok{();}
\NormalTok{    cout }\OperatorTok{\textless{}\textless{}} \StringTok{"md5:"} \OperatorTok{\textless{}\textless{}}\NormalTok{ hash }\OperatorTok{\textless{}\textless{}}\NormalTok{ endl}\OperatorTok{;}
    \ControlFlowTok{return}\NormalTok{ hash}\OperatorTok{.}\FunctionTok{substr}\OperatorTok{(}\DecValTok{0}\OperatorTok{,} \DecValTok{4}\OperatorTok{)} \OperatorTok{+}\NormalTok{ salt }\OperatorTok{+}\NormalTok{ hash}\OperatorTok{.}\FunctionTok{substr}\OperatorTok{(}\DecValTok{4}\OperatorTok{);}
\OperatorTok{\}}
\end{Highlighting}
\end{Shaded}

于是我们就有了一个key生成器了，完整的key生成器源码在Github。


\end{document}
